\documentclass[a4paper,12pt]{article}

\usepackage[utf8]{inputenc}
\usepackage[ukrainian]{babel}
\usepackage{amsmath, amssymb}
\usepackage{graphicx}
\usepackage{tcolorbox}
\usepackage{geometry}
\usepackage{fancyhdr} % ДОДАНО

\geometry{margin=2cm}

\newcommand{\student}{Левенець Владислава}
\newcommand{\labdate}{\today}

\pagestyle{fancy} % ДОДАНО
\fancyhead[L]{\student}
\fancyhead[R]{\labdate}

\title{Стереометрія. Призма і паралелепіпед}
\author{}
\date{}

\begin{document}

\maketitle

\section*{Призма}

\begin{tcolorbox}[colback=blue!5!white,colframe=blue!75!black,title=Визначення]
Призма — многогранник, дві грані якого є рівними многокутниками,
які лежать у паралельних площинах (основи призми), а решта граней —
паралелограми (бічні грані).
\end{tcolorbox}

Висота призми — відрізок перпендикуляра, опущеного з точки однієї основи
на площину іншої основи.

\begin{figure}[h]
    \centering
    \includegraphics[scale=0.8]{foto11.png}
    \caption{Зображення призми (прямої та похилої)}
    \label{fig:prism}
\end{figure}

\textbf{Позначення:}
\begin{itemize}
    \item $l$ — бічне ребро;
    \item $P$ — периметр основи;
    \item $S_{\text{осн}}$ — площа основи;
    \item $H$ — висота;
    \item $V$ — об’єм;
    \item $S_{\text{біч}}$ — площа бічної поверхні;
    \item $S_{\text{повн}}$ — площа повної поверхні.
\end{itemize}

\textbf{Формули:}
\[V = S_{\text{осн}} \cdot H\]
\[S_{\text{біч}} = P \cdot H\]
\[S_{\text{повн}} = S_{\text{біч}} + S_{\text{осн}}\]

\section*{Паралелепіпед}

Паралелепіпед призма, основами якої є паралелограми. Усі шість граней паралелепіпеда паралелограми.

Діагоналі паралелепіпеда відрізки, які сполучають вершини паралелепіпеда, що не належать одній і тій же грані.

Паралелепіпед, бічні ребра якого перпендикулярні до площини основи, називається прямим 

Прямий паралелепіпед, основою якого є прямокут-ник, називається прямокутним 

Якщо а, в, с виміри, а діагональ прямокутного паралелепіпеда, то

\textbf{Формули:}
\[d^2 = a^2 + b^2 + c^2\]
\[V = abc\]
\[S_{\text{повн}} = 2(ab + bc + ac)\]

\begin{figure}[h]
    \centering
    \includegraphics[scale=0.5]{foto22.png}
    \caption{(прямокутна призма)}
    \label{fig:prism}
\end{figure}


\end{document}